% 20-06(20쪽) — y=f(x) 의 그래프. x<0 은 기울기 -1 인 직선, 0<=x<1 은 y=-2, 1<x<3 은 (1,-2)에서 (3,0)으로 오르는 곡선, x>3 은 내려가는 직선.
% 원문의 빈 점·찬 점과 점선(1 눈금·-1 눈금 연결)만 옮긴다. 스캔 화질이 낮아 TikZ 로 다시 그림.
\begin{tikzpicture}[line width=0.55pt, x=0.62cm, y=0.62cm, every node/.style={font=\footnotesize}]
  \draw[->, >=stealth, line width=0.7pt] (-2.4,0) -- (5.0,0) node[below=1pt] {$x$};
  \draw[->, >=stealth, line width=0.7pt] (0,-3.0) -- (0,1.9) node[left=1pt] {$y$};
  % 점선
  \draw[densely dashed, line width=0.4pt] (0,-1) -- (1,-1) (1,-2) -- (1,0);
  \draw[densely dashed, line width=0.4pt] (0,1) -- (3,1) (3,0) -- (3,1);
  % x<0 : 기울기 -1
  \draw[line width=0.6pt] (-2.2,1.2) -- (0,-1);
  % 0<=x<1 : y=-2
  \draw[line width=0.6pt] (0,-2) -- (1,-2);
  % 1<x<=3 : 위로 볼록한 곡선
  \draw[line width=0.6pt] (1,-2) .. controls (1.5,-0.85) and (2.2,-0.12) .. (3,0);
  % x>3 : 기울기 -1
  \draw[line width=0.6pt] (3,1) -- (4.7,-0.7);
  % 빈 점
  \draw[fill=white, line width=0.5pt] (0,-1) circle (2.2pt);
  \draw[fill=white, line width=0.5pt] (1,-2) circle (2.2pt);
  \draw[fill=white, line width=0.5pt] (3,1) circle (2.2pt);
  % 찬 점
  \fill (0,-2) circle (2.2pt);
  \fill (1,-1) circle (2.2pt);
  \fill (3,0) circle (2.2pt);
  % 라벨
  \node[anchor=east, inner sep=2.5pt] at (-0.14,1) {$1$};
  \node[anchor=east, inner sep=2.5pt] at (-0.14,-1) {$-1$};
  \node[anchor=east, inner sep=2.5pt] at (-0.14,-2) {$-2$};
  \node[below=1pt] at (-1.1,-0.05) {$-1$};
  \node[below right=-2pt and -3pt] at (0,0) {$\pt{O}$};
  \node[above=2pt] at (1,0.05) {$1$};
  \node[below=1pt] at (3,-0.05) {$3$};
  \node[below=1pt] at (4,-0.05) {$4$};
  \node[anchor=west] at (1.5,1.75) {$y=f(x)$};
\end{tikzpicture}
